\subsection{PINN Reproducibility: A Cross-Cutting Observation}\label{sec:pinn-repro}

Three of the five Wave~4 papers (Eivazi~2022, Kopani\'cakov\'a~2023,
Zhang~2019) share a striking failure pattern: the \emph{methodology}
reproduces qualitatively (correct physics, correct algorithmic structure,
correct trend directions), but \emph{absolute numerical precision} does not.
Paper-claimed errors of 0.01--1\% become 1--50\% in our hands---a gap of
10--100$\times$.

The root causes are consistent across all three:
\begin{enumerate}[leftmargin=1.4em,itemsep=0.2em,topsep=0.2em]
\item \textbf{No public code.} All three papers lack published implementations.
  The two REPLICATED papers in this batch (Kochkov~2021, Grossmann~2023) both
  have public repositories, underscoring the importance of code availability.
\item \textbf{L-BFGS implementation sensitivity.} PINNs rely heavily on L-BFGS
  for the final refinement phase. Line-search strategy (cubic backtracking +
  strong Wolfe vs Armijo vs fixed step), history length, and convergence
  tolerances can shift final accuracy by 1--2 orders of magnitude. Standard
  library L-BFGS (PyTorch, scipy) does not match custom implementations.
\item \textbf{Undocumented hyperparameters.} Loss balancing coefficients,
  collocation point density, learning rate schedules, initialization seeds,
  and curriculum strategies are often omitted or incompletely specified. PINN
  training is known to be highly sensitive to these choices.
\item \textbf{Reference data availability.} Eivazi~2022 requires DNS/LES
  datasets from KTH that are not publicly available. Without the correct
  boundary data, the PINN cannot learn the correct field.
\end{enumerate}

\noindent\textbf{Contrast with the REPLICATED papers.} Kochkov~2021
(jax-cfd) provides open-source code, public training data on GCS, and
pre-trained model outputs---enabling independent verification. Grossmann~2023
provides a complete GitHub repository with all benchmark configurations.
Both achieved REPLICATED status despite being technically demanding. The
dichotomy is stark: \emph{code availability is the single strongest predictor
of replication success for PINN/ML-PDE papers.}

\noindent\textbf{Implication for the field.} The PINN reproducibility gap we
observe is consistent with findings by Krishnapriyan et al.\ (2021,
``Characterizing possible failure modes in physics-informed neural
networks'') and suggests that PINN papers should be held to a higher
standard of implementation disclosure than classical numerical methods,
precisely because their results are more sensitive to implementation
details.


\clearpage

